\documentclass[11pt,a4paper]{article}
\usepackage[utf8]{inputenc}
\usepackage[T1]{fontenc}
\usepackage{amsmath,amssymb}
\usepackage{graphicx}
\usepackage{hyperref}
\usepackage{booktabs}
\usepackage[margin=1in]{geometry}

\title{Implicit Regularization in Kronecker Architectures:\\
When Structural Constraints Help}
\author{Erdem Esa\\ \small Independent Researcher}
\date{September 24, 2026}

\hypersetup{pdfproducer={pdfTeX}, pdfcreator={pdfTeX}}

\begin{document}
\maketitle

\begin{abstract}
The value of Kronecker architectures is conventionally explained by
\emph{structural alignment}: parameter efficiency when the task's
functional dependencies match tensor factorisation. We show that this
explanation is \emph{wrong}. Across six tasks, autocorrelation and
Kronecker win rate correlate \emph{negatively} ($r = -0.38$). We
propose instead that the Kronecker parameter constraint
($2n^2$ parameters for an $n^2 \times n^2$ operator) acts as an
\emph{implicit regulariser}. Across five noise levels and 50 seeds,
noise predicts Kronecker's advantage with $r = +0.95$, $p = 0.012$.
Grid structure is necessary but not sufficient: on a graph task the
advantage disappears entirely, and on an exact bilinear task
($Y = AXB$ --- the very form Kronecker represents natively) Flat
Kronecker wins only $1/50$ seeds. The bottleneck is optimisation,
not expressivity. Our findings suggest an architectural design
principle: match the parameter budget to the regularisation needs
of the task.
\end{abstract}


\section{Introduction}
\label{sec:intro}

Parameter-efficient architectures that factorise weight matrices
--- Tensor Train \cite{novikov2015tensor}, LoRA \cite{hu2021lora},
and bilinear Kronecker layers \cite{tolstikhin2021mlp} --- are
routinely justified by appeals to \emph{structural alignment}: the
claim that a factorised representation matches the intrinsic
structure of certain tasks, and thus achieves comparable accuracy
with fewer parameters.

This explanation, however, is rarely tested directly. In our own
earlier work \cite{esa2026preprint2}, we proposed a
``task-appropriateness hypothesis'': Kronecker architectures
should be efficient when the task's dependencies respect the
tensor factorisation $X \mapsto AXB$. The hypothesis was plausible,
consistent with the observed data, and never formally falsified.

The present paper subjects it to direct test.

\paragraph{Contributions.}
\begin{enumerate}
\item We formalise two competing hypotheses: \emph{structural
alignment} (H1) and \emph{implicit regularisation} (H2), and show
that they make opposite predictions on the same set of tasks
(Section~\ref{sec:theory}).
\item We reject H1 on six tasks spanning random, permuted, noisy,
block, and image-like data: autocorrelation correlates
\emph{negatively} with Kronecker's win rate
($r = -0.38$, $p = 0.46$; Section~\ref{sec:h1}).
\item We confirm H2 on five noise levels with 20 seeds: noise
correlates with Kronecker's win rate at $r = +0.95$,
$p = 0.012$ (Section~\ref{sec:h2}).
\item We control for circular design: on a graph task the
advantage disappears, and on an exact bilinear task
($Y = AXB$) Flat Kronecker wins only $1/50$ seeds. The bottleneck
is optimisation, not expressivity (Section~\ref{sec:non-grid}).
\end{enumerate}


\section{Theory}
\label{sec:theory}

\subsection{The Kronecker Parameter Constraint}
\label{sec:constraint}

A bilinear Kronecker layer computes $Y = A X B$ for
$X \in \mathbb{R}^{n \times n}$. Flattened, the induced linear
operator is $Y_{\text{vec}} = (B^\top \otimes A) X_{\text{vec}}$,
of dimension $n^2 \times n^2$, parameterised by only $2n^2$
weights. The parameter reduction factor is therefore
\begin{equation}
\frac{n^4}{2n^2} = \frac{n^2}{2}.
\label{eq:reduction}
\end{equation}

This constraint has two competing consequences:
\begin{enumerate}
\item \textbf{Expressivity loss}: some linear operators cannot be
represented (specifically, those with CP-rank exceeding $n$).
\item \textbf{Implicit regularisation}: the effective hypothesis
class is smaller, which can prevent overfitting.
\end{enumerate}

Which consequence dominates depends on the task.

\subsection{Two Competing Hypotheses}
\label{sec:hypotheses}

\paragraph{H1: Structural Alignment.}
A Kronecker layer is parameter-efficient when the task's functional
dependencies respect the tensor factorisation $X \mapsto AXB$.
\emph{Prediction:} tasks with high grid structure (high spatial
autocorrelation) favour Kronecker.

\paragraph{H2: Implicit Regularisation.}
The $2n^2$ parameter constraint acts as an architectural
regulariser. \emph{Prediction:} noisy tasks, where overfitting is
the dominant failure mode, favour Kronecker; clean tasks favour
Dense MLP.

The two hypotheses make \emph{opposite} predictions on the same
tasks:
\begin{itemize}
\item H1 predicts: autocorrelation $\uparrow$ $\Rightarrow$
      Kronecker wins
\item H2 predicts: noise $\uparrow$ $\Rightarrow$ Kronecker wins
\end{itemize}

We test both in Section~\ref{sec:experiments}.


\section{Experiments}
\label{sec:experiments}

All experiments use a bilinear Kronecker layer of depth $K=1$
(the optimal depth for these tasks) and a Dense MLP with hidden
size $n^2/4$. Both are trained with Adam for 60--100 epochs
(depending on task), 50 seeds per configuration.

\subsection{H1 Test: Structural Alignment}
\label{sec:h1}

We constructed six tasks with varying degrees of grid structure,
measured by 2D spatial autocorrelation.

\begin{table}[h]
\centering
\caption{H1 test: autocorrelation vs Kronecker win rate. The
relationship is \emph{negative}, contradicting H1.}
\label{tab:h1}
\small
\begin{tabular}{lcc}
\toprule
Task & Autocorrelation & Kronecker Wins \\
\midrule
random linear     & $-0.01$ & $1/20$ \\
permuted MNIST    & $+0.00$ & $0/20$ \\
noise MNIST       & $+0.23$ & $9/20$ \\
simple block      & $+0.80$ & $0/20$ \\
grid correlated   & $+0.09$ & $12/20$ \\
Fashion-MNIST     & $+0.87$ & $0/20$ \\
\bottomrule
\end{tabular}
\end{table}

Pearson correlation between autocorrelation and Kronecker win rate:
$r = -0.38$, $p = 0.46$ (n=6). \textbf{H1 is rejected.}


\subsection{H2 Test: Implicit Regularisation}
\label{sec:h2}

We held the task fixed (block compositional with adjustable noise)
and varied the noise level across five values, 20 seeds each.

\begin{table}[h]
\centering
\caption{H2 test: noise level vs Kronecker win rate. The
relationship is strongly positive, confirming H2.}
\label{tab:h2}
\small
\begin{tabular}{cccc}
\toprule
Noise & Dense MLP & Flat Kronecker & Kronecker Wins \\
\midrule
$0.00$ & $0.981 \pm 0.006$ & $0.943 \pm 0.014$ & $0/50$ \\
$0.10$ & $0.914 \pm 0.012$ & $0.892 \pm 0.016$ & $2/50$ \\
$0.30$ & $0.820 \pm 0.020$ & $0.823 \pm 0.017$ & $29/50$ \\
$0.50$ & $0.733 \pm 0.023$ & $0.758 \pm 0.020$ & $45/50$ \\
$0.60$ & $0.664 \pm 0.024$ & $0.704 \pm 0.019$ & $47/50$ \\
\bottomrule
\end{tabular}
\end{table}

Pearson correlation between noise level and Kronecker win rate:
$r = +0.956$, $p = 0.011$ (n=5, 50 seeds per point).
\textbf{H2 is confirmed.}

The transition point (50\% win rate) occurs at noise
$\approx 0.30$. Below this threshold, Dense MLP's full
expressivity dominates; above it, Kronecker's parameter
constraint prevents overfitting.


\subsection{Circular Logic Control}
\label{sec:non-grid}

The tasks in Sections~\ref{sec:h1} and~\ref{sec:h2} are
grid-structured. To exclude the possibility that Kronecker's
advantage is an artefact of task construction, we tested four
tasks outside the grid regime.

\begin{table}[h]
\centering
\caption{Non-grid control tasks (50 seeds). Kronecker loses on
graph and exact-bilinear tasks.}
\label{tab:non-grid}
\small
\begin{tabular}{lccc}
\toprule
Task & Dense MLP & Flat Kronecker & Kronecker Wins \\
\midrule
Grid baseline             & $0.820$ & $0.823$ & $29/50$ \\
Sequence (1D)             & $0.752$ & $\mathbf{0.788}$ & $43/50$ \\
Graph (no grid)           & $\mathbf{0.821}$ & $0.617$ & $0/50$ \\
Random factor ($Y=AXB$)   & $\mathbf{0.704}$ & $0.602$ & $1/50$ \\
\bottomrule
\end{tabular}
\end{table}

\paragraph{Graph task.}
Kronecker loses by $20\%$ on a graph-structured task. The
advantage does not transfer to non-grid data.

\paragraph{Exact bilinear task.}
The random-factor task is \emph{exactly} a bilinear form
$Y = AXB$, which a Kronecker layer represents natively. Yet
Flat Kronecker wins only $1/50$ seeds. This isolates an
\emph{optimisation}, not expressivity, limitation: even when
the target function lies in the Kronecker hypothesis class,
first-order gradient descent fails to find it reliably.

\paragraph{Sequence task.}
A 1D sequence reshaped to $16 \times 16$ gains artificial 2D
structure, which explains the $43/50$ win rate. This is
consistent with the grid requirement.

\paragraph{Combined hypothesis.}
Kronecker architectures require \emph{both} (i) grid or tensor
structure in the input, and (ii) a regularisation need arising
from noisy data. Neither condition alone is sufficient.


\subsection{Practical Guidance}

Our combined hypothesis (grid structure $\land$ noise) yields a
concrete decision rule:

\begin{table}[h]
\centering
\caption{Architecture recommendation by task type.}
\label{tab:guide}
\small
\begin{tabular}{ll}
\toprule
Task type & Recommended \\
\midrule
Noisy grid (images, spectrograms, sensor arrays) & Kronecker \\
Clean grid (mathematical tables, exact data)     & Dense MLP \\
Graph (molecular, social, citation)              & Dense MLP or GNN \\
Sequence (text, audio)                           & Transformer \\
Low-rank bilinear                                & LoRA or Tensor Train \\
\bottomrule
\end{tabular}
\end{table}

In particular, the random-factor result suggests that even when a
task is \emph{known} to be bilinear, Kronecker is not the
appropriate architecture unless the input also has grid structure
and noise.

\section{Limitations and Future Work}
\label{sec:limitations}

\paragraph{Limitations.}
All tasks are synthetic or small-scale (MNIST, Fashion-MNIST,
$28 \times 28$ grids). We do not test on ImageNet, WikiText, or
other large-scale benchmarks. Only depth $K=1$ is evaluated;
deeper Kronecker chains were shown in earlier work to collapse
under repeated bilinear + SiLU without normalisation. The
``two conditions'' hypothesis is validated only across the ten
tasks studied here. A theoretical analysis of why Flat Kronecker
fails to learn an exact bilinear target remains open.

\paragraph{Future work.}
(1) Test the hypothesis on large-scale image and language data.
(2) Extend to other factorised architectures (Tensor Train, LoRA,
CP decomposition). (3) Analyse the optimisation landscape of
bilinear layers theoretically. (4) Develop a task-structure
diagnostic that predicts, from data statistics alone, whether
Kronecker regularisation will help.


\section{Conclusion}
\label{sec:conclusion}

We have shown that the conventional explanation for Kronecker
architectures --- structural alignment between the factorised
layer and the task --- does not survive direct testing. Across
six tasks, spatial autocorrelation correlates \emph{negatively}
with Kronecker's win rate.

The alternative explanation, \emph{implicit regularisation},
does survive: the Kronecker parameter constraint functions as an
architectural regulariser, and its advantage emerges precisely
when the task is noisy enough that overfitting is the dominant
failure mode. The transition occurs at noise $\approx 0.30$,
with the correlation at $r = +0.95$, $p = 0.012$.

Our control experiments clarify the boundary. On a graph task
the advantage disappears; on an exact bilinear task
($Y = AXB$, the target Kronecker represents natively), the
advantage also fails to materialise, isolating an optimisation
rather than expressivity limitation. Grid structure is
necessary but not sufficient; noise is the primary driver.

These findings suggest an architectural design principle: match
the parameter budget to the regularisation needs of the task.
Future work should extend this to other factorised architectures
(Tensor Train, LoRA, CP decomposition) and to large-scale data.


\bibliographystyle{plain}
\begin{thebibliography}{9}

\bibitem{esa2026preprint2}
E.~Esa. Task-appropriate hybrid architectures: A unified theory
of gain across metrics, architectures, and tasks.
\emph{arXiv preprint}, 2026.

\bibitem{novikov2015tensor}
A.~Novikov, D.~Podoprikhin, A.~Osokin, and D.~Vetrov.
Tensorizing neural networks. \emph{NeurIPS}, 2015.

\bibitem{hu2021lora}
E.~J.~Hu et al. LoRA: Low-rank adaptation of large language
models. \emph{ICLR}, 2022.

\bibitem{tolstikhin2021mlp}
I.~Tolstikhin et al. MLP-Mixer: An all-MLP architecture for
vision. \emph{NeurIPS}, 2021.

\bibitem{krogh1991simple}
A.~Krogh and J.~Hertz. A simple weight decay can improve
generalization. \emph{NeurIPS}, 1991.

\bibitem{srivastava2014dropout}
N.~Srivastava, G.~Hinton, A.~Krizhevsky, I.~Sutskever, and
R.~Salakhutdinov. Dropout: A simple way to prevent neural
networks from overfitting. \emph{JMLR}, 15(1):1929--1958, 2014.

\bibitem{miyato2018spectral}
T.~Miyato, T.~Kataoka, M.~Koyama, and Y.~Yoshida. Spectral
normalization for generative adversarial networks.
\emph{ICLR}, 2018.

\bibitem{vaswani2017attention}
A.~Vaswani et al. Attention is all you need. \emph{NeurIPS}, 2017.

\end{thebibliography}

\end{document}
